% ============================================================
%  附录A 推迟格林函数（原书印刷页 533–539）
%  底本：高等教育出版社《固体理论》第二版
%  （PDF 页 = 印刷页 + 20，原书 PDF 页 553–559）
%  公式编号自检：全附录 (A.1)–(A.32) 连续 32 个，无缺号。
%    任务单 tags=31 系 (A.23) 的编号嵌于 OCR 乱码未被计入
%    （原书 p537 在号）；另 (A.12)(A.15)(A.16)(A.21)(A.31) 的
%    OCR tag 残缺（tag{2}/tag{..15}/tag{16}/tag{21}/tag{.31}），
%    均已对照原书页面归位。
%
%  已修正的 OCR/原书错误（除全角句点转"."、删 ~ 间隔符、
%  拆开行内公式与中文等常规项外）：
%  1. 特性 1 小标题 OCR 乱码（"iħℤℋ G_r f{±…}"），按 p533 复原
%     为"1. G_r(t,t')=G_r(t-t')，说明 G_r 代表传播子"。
%  2. 脚注①"可参考祖巴烈夫的总结性文章 Zubarev D N. Soviet
%     Phys. Usp., 1960, 3: 320"OCR 漏失，按 p533 页脚补录。
%  3. (A.4) 各指数因子 OCR 的 h 全部为 ħ 之误识，全附录统一 ħ。
%  4. FM 基态平均式（(A.8) 之下第二例）原书 p534 印
%     "Ŝ_l^+(t)(t)Ŝ_{l'}^-(0)"，(t) 系原书衍字，删一。
%  5. (A.12)(A.15)(A.16) 三块 OCR 乱码按 p535–536 复原；
%     (A.12) 第二行行末 ± 号后接 <0|B|n>…（照原书断行）。
%  6. (A.17) 分式结构按 p536 复原：{exp[β(E_n−E_m)]±1}/
%     {ħω−(E_n−E_m)+iη}。
%  7. (A.18) OCR"e^{βħω_m}"按 p536 复原为 e^{βħω_{nm}}；
%     "\hat{h}ω"、"\hat{⊗}"分别复原为 ħω、δ。
%  8. (A.21) 第三矩阵元 OCR"<m|A|m>"按 p537 复原为 <m|A|n>。
%  9. (A.23) 整块 OCR 乱码按 p537 复原（含右端按上半/下半复平面
%     取 G_r/G_a 的分段表示）。
% 10. (A.26) 分子 [exp(βE'_{nm})±1] 按原书补全（OCR 括号残缺）。
% 11. 原书 p538 (A.29) 印"(e^{βE'_{nm}}∓1)"，按 (A.26) 分子及
%     代回 (A.28) 应恰得 <BA> 判定 ∓ 为 ± 之误，改 ±1。
% 12. (A.30)(A.31) 指数因子 OCR 均作 e^{-βE'_n}±e^{-βE'_n}，按
%     p539 复原为 e^{-βE'_m}±e^{-βE'_n}。
% 13. 原书 p539 末段"我们还将用它建立利用它建立介电函数的严格
%     公式"衍"利用它建立"四字（与上下文重复），删。
% 14. 体例：< <A|B> > 统一 \langle\langle…\rangle\rangle；花体
%     ℋ、ℰ 统一 \mathscr；虚单位 i、指数底 e 及 Tr 用正体。
%  遗留问题：
%  · 无。
% ============================================================
\chapter{推迟格林函数}

在海森伯绘景中任意两个算符 $A$ 与 $B$ 所组成的函数
\begin{equation}
\begin{aligned}
G_{\mathrm{r}}(t,t')&=-\frac{\mathrm{i}}{\hbar}\theta(t-t')
\langle[A(t),B(t')]_{\pm}\rangle\\
&=\left\{
\begin{aligned}
&-\frac{\mathrm{i}}{\hbar}\bigl(\langle A(t)B(t')\rangle
\pm\langle B(t')A(t)\rangle\bigr)\qquad(t>t')\\
&0\qquad(t<t')
\end{aligned}\right.\\
&\equiv\langle\langle A(t);B(t')\rangle\rangle_{\mathrm{r}}\text{，}
\end{aligned}
\end{equation}
称为双时推迟格林函数. $G_{\mathrm{r}}$ 是玻戈留玻夫等人首先引入的\footnote{可参考祖巴烈夫的总结性文章 Zubarev D N. Soviet Phys. Usp., 1960, 3: 320}，$\langle\langle A(t);B(t')\rangle\rangle_{\mathrm{r}}$ 是双时推迟格林函数的玻氏记号，其中泊松括号下角的 $\pm$ 号可根据方便与需要选取. 当 $A$ 和 $B$ 属于费密型算子时，满足反对易律，取 $+$ 号比较方便，而对于玻色型算子，满足对易律，往往取 $-$ 号. 式(A.1)中 $\langle\cdots\rangle$ 代表统计平均，对于正则系综
\begin{equation}
\langle A\rangle=Z^{-1}\mathrm{Tr}(\mathrm{e}^{-\beta H}A)
\text{，}\quad Z=\mathrm{Tr}(\mathrm{e}^{-\beta H})
\text{，}\quad\beta=(k_{\mathrm{B}}T)^{-1}
\end{equation}
$G_{\mathrm{r}}$ 是实时间的温度格林函数，代表海森伯绘景中的一种有用的时间关联函数，有下列重要特性：

\subsection*{1. $G_{\mathrm{r}}(t,t')=G_{\mathrm{r}}(t-t')$，说明 $G_{\mathrm{r}}$ 代表传播子}

由于算符的乘积在 $\mathrm{Tr}$ 号内具有循环性
\begin{equation}
\mathrm{Tr}(ABC)=\mathrm{Tr}(CAB)=\mathrm{Tr}(BCA)
\end{equation}
因此
\begin{equation}
\begin{aligned}
\langle A(t)B(t')\rangle&=Z^{-1}\mathrm{Tr}
\bigl(\mathrm{e}^{-\beta H}\mathrm{e}^{\frac{\mathrm{i}}{\hbar}Ht}
A\mathrm{e}^{-\frac{\mathrm{i}}{\hbar}H(t-t')}
B\mathrm{e}^{-\frac{\mathrm{i}}{\hbar}Ht'}\bigr)\\
&=Z^{-1}\mathrm{Tr}\bigl(\mathrm{e}^{-\beta H}
\mathrm{e}^{\frac{\mathrm{i}}{\hbar}H(t-t')}
A\mathrm{e}^{-\frac{\mathrm{i}}{\hbar}H(t-t')}B\bigr)\\
&=Z^{-1}\mathrm{Tr}\bigl[\mathrm{e}^{-\beta H}A(t-t')B\bigr]\\
&=\langle A(t-t')B(0)\rangle
\end{aligned}
\end{equation}
同理可证明
\begin{equation}
\langle B(t')A(t)\rangle=\langle B(0)A(t-t')\rangle
\end{equation}
于是双时推迟格林函数
\begin{equation}
\begin{aligned}
G_{\mathrm{r}}(t-t')&=-\frac{\mathrm{i}}{\hbar}\theta(t-t')
\langle[A(t),B(t')]_{\pm}\rangle\\
&=-\frac{\mathrm{i}}{\hbar}\theta(t-t')
\langle[A(t-t'),B(0)]_{\pm}\rangle\\
&=G_{\mathrm{r}}(t-t')
\end{aligned}
\end{equation}
说明 $G_{\mathrm{r}}$ 只是时间差 $(t-t')$ 的函数，特性 1 得证. 根据式(A.6)，可以简单的取 $t'=0$，并令 $t$ 代表时间差，则推迟格林函数的定义简化为
\begin{equation}
G_{\mathrm{r}}(t)\equiv\langle\langle A(t);B\rangle\rangle_{\mathrm{r}}
=-\frac{\mathrm{i}}{\hbar}\theta(t)\langle[A(t),B]_{\pm}\rangle
\end{equation}
其中 $B=B(0)$，而
\begin{equation}
A(t)=\mathrm{e}^{\frac{\mathrm{i}}{\hbar}Ht}
A\mathrm{e}^{-\frac{\mathrm{i}}{\hbar}Ht}
\end{equation}
$G_{\mathrm{r}}(t)$ 与实验测量直接有关，是最重要的一种格林函数. 例如，设 $A=\hat{S}_{l}^{+}$ 和 $B=\hat{S}_{l'}^{-}$ 所组成的双时推迟格林函数在 $t>0$ 时为
\begin{equation*}
\begin{aligned}
G_{\mathrm{r}}(t)&=\langle\langle\hat{S}_{l}^{+}(t);
\hat{S}_{l'}^{-}(0)\rangle\rangle_{\mathrm{r}}\\
&=-\frac{\mathrm{i}}{\hbar}\bigl\{\langle\hat{S}_{l}^{+}(t)
\hat{S}_{l'}^{-}(0)\rangle
-\langle\hat{S}_{l'}^{-}(0)\hat{S}_{l}^{+}(t)\rangle\bigr\}
\end{aligned}
\end{equation*}
当 $T=0$ K 时，对于铁磁自旋晶格 $G_{\mathrm{r}}(t)$ 由 FM 相的基态平均表示
\begin{equation*}
\begin{aligned}
G_{\mathrm{r}}(t)&=-\frac{\mathrm{i}}{\hbar}\bigl\{\langle0|
\hat{S}_{l}^{+}(t)\hat{S}_{l'}^{-}(0)|0\rangle
\pm\langle0|\hat{S}_{l'}^{-}(0)\hat{S}_{l}^{+}(t)|0\rangle\bigr\}\\
&=-\frac{\mathrm{i}}{\hbar}\langle0|\hat{S}_{l}^{+}(t)
\hat{S}_{l'}^{-}(0)|0\rangle
\end{aligned}
\end{equation*}
显然，$G_{\mathrm{r}}(t)$ 代表 $t=0$ 时刻在 $l'$ 格点上产生自旋偏离，经过 $t$ 时间在 $l$ 格点所测得的信息. 因此，$\langle\langle\hat{S}_{l}^{+}(t);\hat{S}_{l'}^{-}(0)\rangle\rangle$ 描述晶格中自旋波的传播过程，说明 $G_{\mathrm{r}}(t)$ 是传播子（propagator）.

$G_{\mathrm{r}}(t)$ 的傅里叶变换为
\begin{equation}
G_{\mathrm{r}}(t)=\frac{1}{2\pi}\int_{-\infty}^{\infty}
\mathrm{d}\omega\,G_{\mathrm{r}}(\omega)
\exp[-\mathrm{i}(\omega+\mathrm{i}\eta)t]\qquad(\eta=+0)
\end{equation}
其逆变换
\begin{equation}
G_{\mathrm{r}}(\omega)=\int_{-\infty}^{\infty}
\mathrm{d}t\,G_{\mathrm{r}}(t)\exp[\mathrm{i}(\omega+\mathrm{i}\eta)t]
\end{equation}
在实际运算时，更多的是讨论格林函数 $G_{\mathrm{r}}(\omega)$ 的特性

\subsection*{2. $G_{\mathrm{r}}(\omega)$ 的极点确定互作用系统的元激发}

为简单起见先讨论 $T=0$ K 时 $G_{\mathrm{r}}(\omega)$ 的特性. 设 $|n\rangle$ 与 $E_{n}$ 代表互作用系统哈密顿量 $H$ 的严格本征态和能量本征值，$|0\rangle$ 和 $E_{0}$ 为 $H$ 的基态和基态能. 利用矩阵元关系
\begin{equation}
\langle0|A(t)|n\rangle=\langle0|A|n\rangle
\exp\biggl\{-\frac{\mathrm{i}}{\hbar}(E_{n}-E_{0})\biggr\}
\end{equation}
并考虑到 $T=0$ K 时式(A.7)中只需取基态平均，可求出
\begin{equation}
\begin{aligned}
G_{\mathrm{r}}(t)&=-\frac{\mathrm{i}}{\hbar}\theta(t)
\bigl\{\langle0|A(t)B|0\rangle\pm\langle0|BA(t)|0\rangle\bigr\}\\
&=-\frac{\mathrm{i}}{\hbar}\theta(t)\sum_{n}\bigl\{
\langle0|A|n\rangle\langle n|B|0\rangle
\exp(-\mathrm{i}\omega_{n0}t)\pm\\
&\qquad\langle0|B|n\rangle\langle n|A|0\rangle
\exp(\mathrm{i}\omega_{n0}t)\bigr\}
\end{aligned}
\end{equation}
其中
\begin{equation}
\omega_{n0}\equiv\frac{1}{\hbar}(E_{n}-E_{0})
\end{equation}
根据式(A.10)，对 $G_{\mathrm{r}}(t)$ 的表示式(A.12)作傅里叶变换
\begin{equation}
\begin{aligned}
G_{\mathrm{r}}(\omega)&=\frac{1}{\hbar}\sum_{n}\biggl\{
\frac{\langle0|A|n\rangle\langle n|B|0\rangle}
{\omega-\omega_{n0}+\mathrm{i}\eta}
\pm\frac{\langle0|B|n\rangle\langle n|A|0\rangle}
{\omega+\omega_{n0}+\mathrm{i}\eta}\biggr\}\\
&\equiv\langle\langle A|B\rangle\rangle_{\omega+\mathrm{i}\eta}
\end{aligned}
\end{equation}
上式称为推迟格林函数的莱曼（Lehmann）表示. 从式(A.14)可知，$G_{\mathrm{r}}(\omega)$ 的极点在下半复平面，它在上半复平面是解析函数，玻戈留玻夫采用记号 $\langle\langle A|B\rangle\rangle_{\omega+\mathrm{i}\eta}$ 明确表示 $G_{\mathrm{r}}(\omega)$ 的这一解析性. 由于 $G_{\mathrm{r}}(\omega)$ 极点的实轴坐标为 $\pm\omega_{n0}$，因此只要能求出格林函数 $G_{\mathrm{r}}(\omega)$，就可以由其极点确定 $\omega_{n0}$，从式(A.13)知道，$\hbar\omega_{n0}$ 是互作用系统激发态 $E_{n}$ 与基态 $E_{0}$ 之差，它正好代表系统的元激发能量. 这就证明了，从格林函数 $G_{\mathrm{r}}(\omega)$ 之极点可决定多体系统的元激发谱. 另一方面，从式(4.5.21)可知，$G_{\mathrm{r}}(\omega)$ 是响应随外场振荡的振幅
\begin{equation*}
\Delta A=G_{\mathrm{r}}(\omega)\exp(-\mathrm{i}\omega t+\eta t)
\end{equation*}
当外扰动频率 $\omega$ 与系统的本征频率 $\omega_{n0}$ 相等时，将发生共振，此时振幅 $G_{\mathrm{r}}(\omega)$ 必为无限大. 所以，从物理分析也说明，格林函数 $G_{\mathrm{r}}(\omega)$ 的极点应给出元激发的信息.

对于有限温度，由于各种本征态 $|n\rangle$ 均以概率 $Z^{-1}\mathrm{e}^{-\beta E_{n}}$ 出现，这时式(A.7)中 $\langle(\cdots)\rangle$ 不再是基态平均 $\langle0|(\cdots)|0\rangle$，而应当取统计平均 $Z^{-1}\sum_{n}\mathrm{e}^{-\beta E_{n}}\langle n|(\cdots)|n\rangle$. 按此，将式(A.12)推广到 $T>0$ K 时应写成
\begin{equation}
\begin{aligned}
G_{\mathrm{r}}(t)&=-\frac{\mathrm{i}}{\hbar}\theta(t)Z^{-1}
\sum_{n,m}\mathrm{e}^{-\beta E_{n}}\bigl\{\langle n|A|m\rangle
\langle m|B|n\rangle\mathrm{e}^{\mathrm{i}\omega_{mn}t}\pm\\
&\qquad\langle n|B|m\rangle\langle m|A|n\rangle
\mathrm{e}^{-\mathrm{i}\omega_{mn}t}\bigr\}
\end{aligned}
\end{equation}
其中 $\hbar\omega_{nm}\equiv(E_{n}-E_{m})$. 再对上式右边第一项作变量代换 $n\leftrightarrow m$，立刻得到 $T>0$ K 时的 $G_{\mathrm{r}}(t)$ 为
\begin{equation}
\begin{aligned}
G_{\mathrm{r}}(t)&=-\frac{\mathrm{i}}{\hbar}\theta(t)Z^{-1}
\sum_{n,m}\exp(-\beta E_{n})\langle n|B|m\rangle
\langle m|A|n\rangle\times\\
&\quad\exp\biggl\{-\frac{\mathrm{i}}{\hbar}(E_{n}-E_{m})t\biggr\}
\bigl\{\exp[\beta(E_{n}-E_{m})]\pm1\bigr\}
\end{aligned}
\end{equation}
经傅里叶变换求得 $T>0$ K 时的莱曼表示
\begin{equation}
\begin{aligned}
G_{\mathrm{r}}(\omega)&=Z^{-1}\sum_{n,m}\exp(-\beta E_{n})
\langle n|B|m\rangle\langle m|A|n\rangle\times\\
&\quad\frac{\exp[\beta(E_{n}-E_{m})]\pm1}
{\hbar\omega-(E_{n}-E_{m})+\mathrm{i}\eta}
\equiv\langle\langle A|B\rangle\rangle_{\omega+\mathrm{i}\eta}
\end{aligned}
\end{equation}
若引入谱函数
\begin{equation}
\begin{aligned}
A(\omega)&\equiv-\frac{1}{\pi}\mathrm{Im}G_{\mathrm{r}}(\omega)\\
&=Z^{-1}\sum_{n,m}\mathrm{e}^{-\beta E_{n}}\langle n|B|m\rangle
\langle m|A|n\rangle\bigl(\mathrm{e}^{\beta\hbar\omega_{nm}}\pm1\bigr)
\delta\bigl[\hbar\omega-(E_{n}-E_{m})\bigr]\\
&=Z^{-1}\sum_{n,m}\langle n|B|m\rangle\langle m|A|n\rangle
\bigl(\mathrm{e}^{-\beta E_{m}}\pm\mathrm{e}^{-\beta E_{n}}\bigr)
\delta\bigl[\hbar\omega-(E_{n}-E_{m})\bigr]
\end{aligned}
\end{equation}
还可导出格林函数 $\langle\langle A|B\rangle\rangle_{\omega}$ 与平均量 $\langle BA\rangle$ 的严格关系——谱定理.

\subsection*{3. 谱定理}

根据式(A.18)，容易证明谱函数 $A(\omega)$ 与分布函数的乘积对频率的积分恰好等于平均量 $\langle BA\rangle$.
\begin{equation}
\begin{aligned}
\int_{-\infty}^{\infty}\frac{A(\omega)}
{\mathrm{e}^{\beta\hbar\omega}\pm1}\mathrm{d}\omega
&=Z^{-1}\sum_{n,m}\mathrm{e}^{-\beta E_{n}}\langle n|B|m\rangle
\langle m|A|n\rangle\\
&=Z^{-1}\sum_{n}\mathrm{e}^{-\beta E_{n}}\langle n|BA|n\rangle\\
&=\langle BA\rangle
\end{aligned}
\end{equation}
再计入 $A(\omega)$ 的定义式(A.18)，求得平均量 $\langle BA\rangle$ 与格林函数 $G_{\mathrm{r}}(\omega)$ 的一个重要关系式
\begin{equation}
\langle BA\rangle=\frac{1}{2\pi}\int_{-\infty}^{\infty}
\frac{\bigl\{-2\,\mathrm{Im}G_{\mathrm{r}}(\omega)\bigr\}}
{\mathrm{e}^{\beta\hbar\omega}\pm1}\mathrm{d}\omega
\end{equation}
上式称为格林函数的谱定理. 在文献中谱定理常被称为涨落-耗散定理（FDT），这是因为 $\mathrm{Im}G_{\mathrm{r}}(\omega)$ 与系统的阻尼有关，而 $\langle BA\rangle$ 可作为描述涨落的关联函数.

推迟格林函数 $\langle\langle A|B\rangle\rangle_{\omega+\mathrm{i}\eta}$ 的极点在下半复平面，它只是上半复平面的解析函数，若将式(A.17)中 $\mathrm{i}\eta$ 换成 $-\mathrm{i}\eta$，则定义一个在下半复平面的解析函数
\begin{equation}
\begin{aligned}
\langle\langle A|B\rangle\rangle_{\omega-\mathrm{i}\eta}
&\equiv Z^{-1}\sum_{n,m}\exp(-\beta E_{n})\langle n|B|m\rangle
\langle m|A|n\rangle\times\\
&\quad\frac{\exp[\beta(E_{n}-E_{m})]\pm1}
{\hbar\omega-(E_{n}-E_{m})-\mathrm{i}\eta}=G_{\mathrm{a}}(\omega)
\end{aligned}
\end{equation}
其极点在下半面，再对 $G_{\mathrm{a}}(\omega)$ 作傅里叶变换
\begin{equation}
\begin{aligned}
G_{\mathrm{a}}(t)&=\frac{1}{2\pi}\int_{-\infty}^{\infty}
\mathrm{d}\omega\,G_{\mathrm{a}}(\omega)
\exp[-\mathrm{i}(\omega-\mathrm{i}\eta)t]\qquad(\eta=+0)\\
&=\frac{\mathrm{i}}{\hbar}\theta(-t)\langle[A(t),B]_{\pm}\rangle
\equiv\langle\langle A(t);B\rangle\rangle_{\mathrm{a}}\\
&=\left\{
\begin{aligned}
&0\qquad(t>0)\\
&\frac{\mathrm{i}}{\hbar}\bigl\{\langle A(t)B\rangle
\pm\langle BA(t)\rangle\bigr\}\qquad(t<0)
\end{aligned}\right.
\end{aligned}
\end{equation}
所求出的时间关联函数逆时序传播，代表超前格林函数，因此用 $G_{\mathrm{a}}(t)$ 与 $G_{\mathrm{a}}(\omega)$ 表示.

利用 $G_{\mathrm{r}}(\omega)$ 与 $G_{\mathrm{a}}(\omega)$ 可定义包括上、下半复平面在内的一个全格林函数
\begin{equation}
\begin{aligned}
G(\omega)&\equiv\langle\langle A|B\rangle\rangle_{\omega}\\
&=Z^{-1}\sum_{n,m}\mathrm{e}^{-\beta E_{n}}\langle n|B|m\rangle
\langle m|A|n\rangle
\frac{\mathrm{e}^{\beta(E_{n}-E_{m})}\pm1}
{\hbar\omega-(E_{n}-E_{m})}\\
&=\left\{
\begin{aligned}
&G_{\mathrm{r}}(\omega)\qquad(\text{上半复平面})\\
&G_{\mathrm{a}}(\omega)\qquad(\text{下半复平面})
\end{aligned}\right.
\end{aligned}
\end{equation}
其中 $\omega$ 为复量. 尽管全格林函数 $G(\omega)$ 在上半面由 $G_{\mathrm{r}}(\omega)$ 定义，而下半面由 $G_{\mathrm{a}}(\omega)$ 定义，但在实轴上却无定义. 这是因为当 $\eta\to+0$ 时，$G_{\mathrm{r}}(\omega)$ 与 $G_{\mathrm{a}}(\omega)$ 分别从上、下岸趋向实轴，它们的实部相等，而虚部则异号，导致在实轴上 $G(\omega)$ 无法定义，人们必须沿实轴作割线才能保证 $G(\omega)$ 的解析性. 但是利用此割线两边 $\mathrm{Im}\,G_{\mathrm{r}}(\omega)$ 与 $\mathrm{Im}\,G_{\mathrm{a}}(\omega)$ 异号和 $\mathrm{Re}G_{\mathrm{r}}(\omega)=\mathrm{Re}G_{\mathrm{a}}(\omega)$ 特征，又可得等式
\begin{equation*}
-2\,\mathrm{Im}G_{\mathrm{r}}(\omega)=\mathrm{i}\bigl\{
\langle\langle A|B\rangle\rangle_{\omega+\mathrm{i}\eta}
-\langle\langle A|B\rangle\rangle_{\omega-\mathrm{i}\eta}\bigr\}
\end{equation*}
并将谱定理改写为
\begin{equation}
\langle BA\rangle=\frac{\mathrm{i}}{2\pi}\int_{-\infty}^{\infty}
\frac{\langle\langle A|B\rangle\rangle_{\omega+\mathrm{i}\eta}
-\langle\langle A|B\rangle\rangle_{\omega-\mathrm{i}\eta}}
{\mathrm{e}^{\beta\hbar\omega}\pm1}\mathrm{d}\omega
\end{equation}
式(A.24)与式(A.20)是等效的，前者是祖巴烈夫最早引入谱定理时的表示式，这两个公式都是由格林函数计算平均量的有用工具.

以上讨论的谱定理和谱表示仅限于正则系综，对于互作用系统，往往需要采用巨正则系综作统计平均
\begin{equation}
\langle X\rangle=Z_{\mathrm{G}}^{-1}\mathrm{Tr}
(\mathrm{e}^{-\beta\mathscr{H}}X)\text{，}\quad
Z_{\mathrm{G}}=\mathrm{Tr}(\mathrm{e}^{-\beta\mathscr{H}})
\text{，}\quad\mathscr{H}=H-\mu\hat{N}
\end{equation}
其中 $Z_{\mathrm{G}}$ 是巨配分函数，$\mu$ 为化学势，$\hat{N}$ 代表总粒子数算符，并取 $\hbar=1$ 单位. 类似于上述讨论，不难证明推迟与超前格林函数在巨正则系综内的莱曼表示式为
\begin{equation}
\begin{aligned}
\langle\langle A|B\rangle\rangle_{\omega\pm\mathrm{i}\eta}
&=Z_{\mathrm{G}}^{-1}\sum_{n,m}\exp(-\beta E_{n}')
\langle n|B|m\rangle\\
&\qquad\times\langle m|A|n\rangle
\frac{\bigl[\exp(\beta E_{nm}')\pm1\bigr]}
{\omega-E_{nm}\pm\mathrm{i}\eta}
\end{aligned}
\end{equation}
其中 $\beta=(k_{\mathrm{B}}T)^{-1}$，$E_{nm}=E_{n}-E_{m}$，$E_{n}$ 是 $H$ 的能量本征值，因此 $E_{nm}$ 代表从 $m$ 态到 $n$ 态的激发能，它正好等于格林函数 $\langle\langle A|B\rangle\rangle_{\omega}$ 的极点，并决定系统的元激发. 上式中的 $E_{n}'$ 代表 $H-\mu\hat{N}$ 的本征值，且 $E_{n}'=E_{n}-\mu N_{n}$，$N_{n}$ 为 $n$ 态时系统的总粒子数，因此
\begin{equation}
E_{nm}'\equiv(E_{n}'-E_{m}')=E_{nm}-\mu\delta N_{nm}=E_{nm}-\mu
\end{equation}
其中 $\delta N_{nm}=N_{n}-N_{m}$ 代表从 $m$ 态激发到 $n$ 态系统总粒子数的变化. 在具体的格林函数中往往取 $B$ 为产生算符，例如单电子格林函数 $\langle\langle C_{k\sigma}|C_{k\sigma}^{+}\rangle\rangle_{\omega+\mathrm{i}\eta}$，这时式(A.26)中出现 $\langle n|C_{k\sigma}^{+}|m\rangle$ 矩阵元，它将限定从 $m\to n$ 仅增加一个粒子，$\delta N_{nm}=1$，从而得到式(A.27)的最后一个等式.

利用莱曼表示式(A.26)和式(A.27)，可以立刻导出巨正则系统的格林函数谱定理为
\begin{equation}
\langle BA\rangle=\frac{\mathrm{i}}{2\pi}\int_{-\infty}^{\infty}
\frac{\langle\langle A|B\rangle\rangle_{\omega+\mathrm{i}\eta}
-\langle\langle A|B\rangle\rangle_{\omega-\mathrm{i}\eta}}
{\mathrm{e}^{\beta(\omega-\mu)}\pm1}\mathrm{d}\omega
\end{equation}
因为
\begin{equation}
\begin{aligned}
&\frac{\mathrm{i}}{2\pi}\bigl\{\langle\langle A|B\rangle\rangle_{\omega+\mathrm{i}\eta}
-\langle\langle A|B\rangle\rangle_{\omega-\mathrm{i}\eta}\bigr\}\\
&\quad=Z_{\mathrm{G}}^{-1}\sum_{n,m}\mathrm{e}^{-\beta E_{n}'}
\langle n|B|m\rangle\langle m|A|n\rangle
\bigl(\mathrm{e}^{\beta E_{nm}'}\pm1\bigr)\delta(\omega-E_{nm})
\end{aligned}
\end{equation}
式(A.28)的等号右边等于
\begin{equation*}
\begin{aligned}
&\frac{\mathrm{i}}{2\pi}\int_{-\infty}^{\infty}\mathrm{d}\omega\,
\frac{\bigl\{\langle\langle A|B\rangle\rangle_{\omega+\mathrm{i}\eta}
-\langle\langle A|B\rangle\rangle_{\omega-\mathrm{i}\eta}\bigr\}}
{\mathrm{e}^{\beta(\omega-\mu)}\pm1}\\
&\quad=Z^{-1}\sum_{n,m}\mathrm{e}^{-\beta E_{n}'}
\langle n|B|m\rangle\langle m|A|n\rangle\\
&\quad=Z^{-1}\sum_{n}\mathrm{e}^{-\beta E_{n}'}\langle n|BA|n\rangle\\
&\quad=\langle BA\rangle
\end{aligned}
\end{equation*}
显然，在巨正则系综的谱定理公式中含有化学势 $\mu$，它是否位于能带之中将直接确定系统是金属还是绝缘体.

类似于式(A.18)的讨论，利用式(A.26)中推迟格林函数的虚部可直接写出巨正则系综的谱函数
\begin{equation}
A(\omega)=Z_{\mathrm{G}}^{-1}\sum_{n,m}\langle n|B|m\rangle
\langle m|A|n\rangle\bigl(\mathrm{e}^{-\beta E_{m}'}
\pm\mathrm{e}^{-\beta E_{n}'}\bigr)
\delta(\omega-E_{n}+E_{m})\text{，}
\end{equation}
由此又可导出谱函数的一个求和规则
\begin{equation}
\begin{aligned}
\int_{-\infty}^{\infty}\mathrm{d}\omega\,A(\omega)
&=Z_{\mathrm{G}}^{-1}\sum_{n,m}\langle n|B|m\rangle
\langle m|A|n\rangle\bigl(\mathrm{e}^{-\beta E_{m}'}
\pm\mathrm{e}^{-\beta E_{n}'}\bigr)\\
&=Z_{\mathrm{G}}^{-1}\sum_{n,m}\mathrm{e}^{-\beta E_{n}'}
\bigl(\langle n|A|m\rangle\langle m|B|n\rangle\pm\\
&\qquad\langle n|B|m\rangle\langle m|A|n\rangle\bigr)\\
&=Z_{\mathrm{G}}^{-1}\sum_{n}\mathrm{e}^{-\beta E_{n}'}
\langle n|[A,B]_{\pm}|n\rangle\\
&=\langle[A,B]_{\pm}\rangle
\end{aligned}
\end{equation}
对于 $A=C$ 和 $B=C^{+}$，当 $C^{+}$ 和 $C$ 分别代表单粒子的产生与消灭算符时，利用玻色子（取 $-$ 号）与费密子（取 $+$ 号）的对易关系 $[C,C^{+}]_{\pm}=1$，上述求和规则简化为
\begin{equation}
\int_{-\infty}^{\infty}\mathrm{d}\omega\,A(\omega)=1
\end{equation}
表明谱函数具有某种归一性. 谱函数在互作用电子系统中有重要应用，我们将在附录 D 中作进一步的介绍.

应当提醒读者，本节所建立的莱曼表示是在假定多体系统哈密顿量 $H$（或 $H-\mu\hat{N}$）的严格本征态 $|n\rangle$ 及能量本征值 $E_{n}$（或 $E_{n}'$）为已知的前提下求得的一个严格的格林函数表示式，实际上人们并不知道 $H$ 与 $H-\mu\hat{N}$ 在一般情况下如何严格求解，因此一般不能用莱曼表示具体计算格林函数，它的主要用处在于证明定理和导出严格关系. 例如，在本节中用于证明谱定理及谱函数的求和规则. 在附录 B 中，我们还将用它建立介电函数的严格公式，用于讨论金属的基态能. 至于具体计算多体系统的格林函数，还需在复平面上建立双时格林函数 $\langle\langle A|B\rangle\rangle_{\omega}$ 的运动方程，并求其近似解. 我们将在附录 C 中介绍格林函数运动方程，而在第十一章中作具体的近似计算.
